\section{Zawartość pracy}

Praca składa się z sześciu rozdziałów. We wstępie omówiono motywację do podjęcia tematu, określono cele do zrealizowania oraz opisano strukturę pracy.
Następny rozdział zawiera przegląd zagrożeń bezpieczeństwa związany z dużymi modelami językowymi oraz klasyfikację ataków promptowych. Zostały w nim także szczegółowo opisane metody bezpośrednich ataków promptowych, z uwzględnieniem ataków opartych na optymalizacji, metod manualnych oraz automatyczny red teaming oparty na modelach.
W rozdziale trzecim omówiono metody obrony przed atakmi na duże modele językowe. Przedstawiono w nim klasyfikację zabezpieczeń w zależności od etapu cyklu życia modelu, mianowicie trenowanie, wdrożenie i ewaluacja. Szczególną uwagę poświącono inżynierii aktywacji. Opisano jej podstawy teoretyczne, mechanizm sterowania modelem w celu zapobiegania atakom oraz również ograniczenia tej techniki.
W podrozdziałach czwartej części pracy przedstawiono metodykę badań. Omówiono cel i założenia eksperymentu, wykorzystane środowisko sprzętowe i oprogramowanie, poddawane ewaluacji modele językowe. Opisano także zbiory danych, wykorzystane metody ataków i techniki mitygacji oraz użyte w badaniu metryki.
Rozdział piąty zawiera analizę wyników badań. Przedstawiono w nim skuteczność ataków bez zastosowania metod mitygacji oraz ewaluację zaimplementowanych metod obrony. Zawiera porównanie skuteczności tych metod, analizę kompromisu między bezpieczeństwem a użytecznością modelu oraz ewaluację wprowadzonych przez mechanizmy obronne opóźnień. Rezultaty są poddane dyskusji, a także sformułowano ograniczenia przeprowadzonych eksperymentów.
Ostatni rozdział stanowi podsumowanie niniejszej pracy magisterskiej. Zawiera końcowe wnioski na bazie otrzymanych rezultatów, a także wskazano w nim możliwości dalszych badań.